\section{Reinforcement Learning from Human Feedback}
\label{sec:rlhf}

The core challenge in deploying large language models is the misalignment between the pretraining objective---predicting the next token on internet text---and the downstream goal of generating outputs that are helpful, honest, and harmless. Reinforcement Learning from Human Feedback (RLHF) emerged as the dominant paradigm for bridging this gap, establishing a systematic pipeline that transforms raw human preference judgments into a training signal capable of steering model behavior. This section traces the intellectual lineage of RLHF, formalizes its three-stage pipeline, and examines the limitations that have motivated subsequent alternatives.

\subsection{Origins and Early Work}

The idea of learning reward functions from human preferences, rather than hand-engineering them, predates the modern LLM era. \cite{christiano2017deep} proposed a framework for \emph{deep reinforcement learning from human preferences}, demonstrating that a reward model trained on pairwise comparisons of trajectory segments could guide an RL agent to solve complex tasks---including Atari games and simulated robot locomotion---without access to the ground-truth reward function. Critically, their approach required human feedback on less than one percent of the agent's total interactions with the environment, making human oversight practically feasible. The reward model was trained by presenting non-expert humans with pairs of short trajectory clips and asking which exhibited more desirable behavior; the resulting learned reward was then optimized via standard policy gradient methods. This work established two foundational principles that would carry forward into the language domain: (i)~pairwise comparisons are a more reliable and cognitively simpler annotation format than absolute scalar ratings, and (ii)~a learned reward model can serve as a proxy objective that decouples the expensive human evaluation from the iterative policy optimization loop.

The transition from simulated environments to natural language was pioneered by \cite{ziegler2019finetuning}, who applied reward learning to fine-tune language models on four tasks: stylistic text continuation (positive sentiment and physically descriptive language) and abstractive summarization on the TL;DR and CNN/DailyMail datasets. Using a pretrained language model as the initial policy, they trained a reward model on approximately 60,000 human comparisons and optimized the policy with Proximal Policy Optimization (PPO)~\cite{schulman2017ppo}. A key technical contribution was the introduction of a Kullback--Leibler (KL) divergence penalty between the learned policy and the original supervised model, which prevented the policy from deviating too far from the pretrained distribution---a regularization mechanism that would become standard in all subsequent RLHF work.

\cite{stiennon2020learning} scaled this approach substantially, applying it to the Reddit TL;DR summarization task with models up to 6.7 billion parameters. They collected a high-quality dataset of 64,832 human comparisons between summaries and demonstrated that policies trained with human feedback significantly outperformed supervised learning baselines---even those with ten times more parameters. Their 6.7B human feedback model produced summaries that human evaluators preferred over the original human-written reference summaries in the dataset. Perhaps most strikingly, the models transferred effectively to out-of-distribution news summarization on CNN/DailyMail without any domain-specific fine-tuning. This work also provided early empirical evidence of \emph{reward over-optimization}: as the policy was optimized more aggressively against the reward model (by reducing the KL penalty coefficient $\beta$), true human preference scores initially improved but eventually degraded, revealing that the reward model's predictions became anti-correlated with genuine quality under heavy optimization.

\subsection{The Three-Stage RLHF Pipeline}

The work of \cite{ouyang2022instructgpt} consolidated and scaled the preceding research into the now-canonical three-stage RLHF pipeline, applying it to align GPT-3 across a broad distribution of real-world user instructions. The resulting system, \textbf{InstructGPT}, became the methodological foundation for ChatGPT and established RLHF as the dominant post-training paradigm for large language models.

\paragraph{Stage 1: Supervised Fine-Tuning (SFT).}
The pipeline begins with supervised fine-tuning of a pretrained language model on a dataset of high-quality demonstrations. In the InstructGPT setting, a team of 40 human contractors wrote demonstrations of desired model behavior for a diverse set of prompts---encompassing generation, open-ended question answering, brainstorming, chat, rewriting, summarization, and classification tasks. A pretrained GPT-3 model was then fine-tuned on approximately 13,000 such demonstration prompts using standard maximum likelihood estimation. The resulting SFT model serves a dual purpose: it provides a strong initialization for the RL policy and establishes a reference distribution against which policy divergence is subsequently penalized.

\paragraph{Stage 2: Reward Model Training.}
The second stage trains a scalar reward model that serves as a proxy for human judgment. Starting from the SFT model with its final unembedding layer replaced by a randomly initialized linear projection head, the reward model is trained to predict which of two candidate completions a human annotator would prefer for a given prompt. \cite{ouyang2022instructgpt} employed a 6B parameter reward model, finding that 175B reward models exhibited training instability.

The reward model is trained under the \textbf{Bradley--Terry} pairwise comparison framework. Given a prompt $x$ and a pair of completions where $y_w$ is preferred over $y_l$ by a human annotator, the probability that $y_w$ is preferred is modeled as:
\begin{equation}
    P(y_w \succ y_l \mid x) = \sigma\!\bigl(r_\theta(x, y_w) - r_\theta(x, y_l)\bigr),
\end{equation}
where $r_\theta(x, y)$ is the scalar reward assigned by the model with parameters $\theta$, and $\sigma(\cdot)$ denotes the logistic sigmoid function. The corresponding training loss is the negative log-likelihood of the observed preferences:
\begin{equation}
    \mathcal{L}_{\text{RM}}(\theta) = -\frac{1}{\binom{K}{2}} \, \mathbb{E}_{(x, y_w, y_l) \sim \mathcal{D}} \Bigl[\log \sigma\!\bigl(r_\theta(x, y_w) - r_\theta(x, y_l)\bigr)\Bigr],
    \label{eq:rm_loss}
\end{equation}
where $\mathcal{D}$ is the dataset of human comparisons and $K$ is the number of candidate responses ranked per prompt. A key efficiency innovation in InstructGPT was presenting labelers with $K \in \{4, \ldots, 9\}$ responses to rank simultaneously, yielding $\binom{K}{2}$ pairwise comparisons per prompt. All comparisons from a single prompt were processed as a single batch element, requiring only one forward pass per completion rather than $\binom{K}{2}$ passes, which both improved computational efficiency and mitigated overfitting caused by the high correlation among comparisons derived from the same ranking task. The reward model outputs were normalized via a bias term so that labeler demonstrations achieved a mean reward of zero before the RL stage.

\paragraph{Stage 3: Reinforcement Learning with PPO.}
In the final stage, the SFT model is fine-tuned as an RL policy to maximize the learned reward signal using Proximal Policy Optimization~\cite{schulman2017ppo}. The environment is formulated as a bandit problem: a prompt $x$ is sampled from the dataset, the policy $\pi_\phi^{\text{RL}}$ generates a response $y$, the reward model assigns a scalar score, and the episode terminates. Each token generation step constitutes an action in the sequential decision process.

To prevent the policy from exploiting imperfections in the reward model---a phenomenon known as \emph{reward hacking}---a per-token KL divergence penalty is added to the reward, constraining the RL policy to remain close to the reference SFT distribution. The augmented objective takes the form:
\begin{equation}
    \mathcal{J}(\phi) = \mathbb{E}_{(x, y) \sim \pi_\phi^{\text{RL}}} \Bigl[r_\theta(x, y) - \beta \, \text{KL}\!\bigl(\pi_\phi^{\text{RL}}(\cdot \mid x) \,\|\, \pi^{\text{SFT}}(\cdot \mid x)\bigr)\Bigr],
    \label{eq:rlhf_objective}
\end{equation}
where $\beta$ controls the strength of the KL penalty. In practice, the KL divergence is computed at the token level as $\sum_t \log \frac{\pi_\phi^{\text{RL}}(y_t \mid x, y_{<t})}{\pi^{\text{SFT}}(y_t \mid x, y_{<t})}$ and subtracted from the reward at each step. The KL penalty serves a dual purpose: it acts as an entropy bonus that encourages exploration and prevents mode collapse, and it constrains the policy to the support of the reward model's training distribution, limiting the scope for reward hacking.

InstructGPT further introduced the \textbf{PPO-ptx} variant, which mixed RL gradients with pretraining gradients on the original language modeling distribution to mitigate performance regressions on standard NLP benchmarks---an instance of the so-called ``alignment tax.'' The combined objective appends a term $\gamma \, \mathbb{E}_{x \sim \mathcal{D}_{\text{pretrain}}}[\log \pi_\phi^{\text{RL}}(x)]$ to Equation~\ref{eq:rlhf_objective}, where $\gamma$ controls the weight of the pretraining loss.

\subsection{InstructGPT: Results and Significance}

The results of \cite{ouyang2022instructgpt} provided compelling evidence for the effectiveness of RLHF at scale. Models were trained at three sizes---1.3B, 6B, and 175B parameters---all using the GPT-3 architecture. The headline finding was striking: \emph{the 1.3B parameter InstructGPT model was preferred by human evaluators over the 175B parameter GPT-3}, despite having more than 100$\times$ fewer parameters. The 175B InstructGPT model was preferred over the 175B GPT-3 baseline 85\% of the time. Beyond raw preference, InstructGPT demonstrated concrete improvements across multiple alignment dimensions. On the TruthfulQA benchmark, InstructGPT generated truthful and informative answers approximately twice as often as GPT-3. Hallucination rates on closed-domain tasks dropped from 41\% to 21\%. Toxic output generation was reduced by approximately 25\% when the model was prompted to be respectful.

These results demonstrated that alignment through human feedback can be far more sample-efficient than scaling model parameters alone. A relatively modest amount of human feedback data---approximately 33,000 comparison prompts for reward model training and 31,000 prompts for PPO---was sufficient to produce dramatic behavioral improvements, suggesting that the \emph{quality} of the training signal matters as much as, or more than, the \emph{quantity} of pretraining compute.

\subsection{Limitations of the Standard RLHF Pipeline}

Despite its empirical success, the standard RLHF pipeline introduces several significant challenges that have motivated subsequent research into simpler alternatives.

\paragraph{Computational Complexity.}
The full RLHF pipeline requires maintaining up to four models simultaneously in GPU memory during the PPO training phase: (i)~the \emph{policy model} being optimized, (ii)~the \emph{reference model} (a frozen copy of the SFT model) for computing the KL penalty, (iii)~the \emph{reward model} for scoring generated completions, and (iv)~the \emph{value/critic model} used by PPO to estimate advantages. For large-scale models, this quadrupling of memory requirements imposes severe infrastructure constraints. \cite{stiennon2020learning} noted the necessity of using completely separate parameters for the policy and value networks to prevent value function updates from destabilizing the pretrained policy, further increasing the memory footprint.

\paragraph{Training Instability.}
PPO-based RLHF training is notoriously sensitive to hyperparameters, including the KL penalty coefficient $\beta$, the PPO clipping ratio, learning rates, and the number of PPO epochs per batch. The interplay between the reward model, the evolving policy, and the value function creates a complex optimization landscape where small perturbations can lead to training collapse or mode collapse. \cite{ouyang2022instructgpt} themselves reported that 175B reward model training was unstable, necessitating the use of a smaller 6B reward model even when training 175B policies.

\paragraph{Reward Hacking and Over-Optimization.}
The reward model is a learned proxy that imperfectly captures the true distribution of human preferences. As the policy is optimized against this proxy, it inevitably discovers outputs that achieve high reward model scores without corresponding improvements in genuine quality---a failure mode known as \emph{reward hacking} or \emph{Goodhart's Law} applied to learned objectives. \cite{stiennon2020learning} provided direct empirical evidence of this phenomenon: as the KL penalty was reduced and the policy optimized more aggressively, reward model scores continued to increase while actual human preference scores first rose and then declined, eventually becoming anti-correlated with the proxy reward. The KL penalty mitigates but does not eliminate this fundamental tension between proxy optimization and true objective alignment.

\paragraph{Human Annotation Bottleneck.}
The pipeline depends critically on the quality, consistency, and scale of human preference annotations. \cite{ouyang2022instructgpt} employed 40 carefully screened contractors, invested in detailed instructions and ongoing calibration, and achieved inter-annotator agreement rates of approximately 73\%. Nevertheless, preference annotations are inherently subjective, expensive to collect at scale, and can encode the biases of the specific annotator pool. The annotation process creates a bottleneck that limits the iteration speed of the overall training pipeline.

\vspace{0.5em}
\noindent These limitations---particularly the computational overhead of maintaining four models and the fragility of PPO-based optimization---have driven the development of alternative approaches. Direct Preference Optimization (DPO) eliminates the reward model entirely by deriving a closed-form policy objective from the same Bradley--Terry preference framework, while Group Relative Policy Optimization (GRPO) replaces the critic model with group-based advantage estimation, reducing the memory footprint while retaining the benefits of RL-based training. These methods, which form the core focus of this work, are discussed in Sections~\ref{sec:dpo} and~\ref{sec:grpo}.
